\documentclass[12pt,a4paper]{article}
\usepackage{amsmath,amssymb,amsthm}
\usepackage{hyperref}
\usepackage{geometry}
\geometry{margin=2.5cm}
\usepackage{enumitem}
\usepackage{booktabs}

\newtheorem{theorem}{Theorem}[section]
\newtheorem{lemma}[theorem]{Lemma}
\newtheorem{corollary}[theorem]{Corollary}
\newtheorem{proposition}[theorem]{Proposition}
\theoremstyle{definition}
\newtheorem{definition}[theorem]{Definition}
\theoremstyle{remark}
\newtheorem{remark}[theorem]{Remark}

\title{The Cosmological Lithium Problem and Recognition Science:\\
$\varphi$-Modulated BBN Reaction Rates}

\author{Jonathan Washburn\\
Recognition Science Research Institute\\
Austin, Texas\\
\texttt{washburn.jonathan@gmail.com}}

\date{March 2026}

\begin{document}
\maketitle

\begin{abstract}
We analyze the cosmological lithium problem within Recognition Science
(RS).  Standard Big Bang Nucleosynthesis (BBN) predicts a primordial
$^7$Li abundance of $^7$Li/H $\approx (4.7 \pm 0.7) \times 10^{-10}$,
while observations of metal-poor halo stars (the Spite plateau) give
$^7$Li/H $\approx (1.6 \pm 0.3) \times 10^{-10}$---a factor of $\sim
3$ discrepancy~\cite{Cyburt2016,Fields2011}.  RS predicts a
$\varphi$-modulated suppression of the key nuclear reaction rate
${}^3\mathrm{He}(\alpha,\gamma){}^7\mathrm{Be}$ via the 8-tick
recognition cycle.  The per-cycle fractional modulation is
$\delta_\varphi = \ln\varphi / 2^D \approx 0.060$ ($6\%$), derived
from the RS $\varphi$-ladder and 8-tick period.  The required factor
of $\sim 3$ suppression is obtained after $N \approx 18$ effective
recognition-cycle epochs; we show that $N \approx 18$ is
order-of-magnitude consistent with the number of $\varphi$-rung
transitions during the ${}^7$Be production window, and note that a
first-principles derivation of $N$ from the RS Hamiltonian remains to
be completed.  The $^4$He abundance and deuterium ratio are insensitive
to this modulation.
\end{abstract}

%% ============================================================
\section{Recognition Science Framework}
\label{sec:framework}
%% ============================================================

The $\varphi$-modulation mechanism applied to BBN lithium production
derives from the unique cost functional fixed by the RS axioms.

\begin{definition}[RS Cost Functional]
For $x > 0$: $J(x) = \tfrac{1}{2}(x + x^{-1}) - 1$.
\end{definition}

\noindent\textbf{Axiom A1 (Normalization):} $J(1) = 0$.\\
\textbf{Axiom A2 (Recognition Composition Law, RCL):}
$J(xy)+J(x/y)=2J(x)J(y)+2J(x)+2J(y)$ for all $x,y>0$.\\
\textbf{Axiom A3 (Calibration):} $J''_{\log}(0) = 1$, where
$J_{\log}(t) := J(e^t)$.

\begin{theorem}[Cost Uniqueness]
\label{thm:unique}
The continuous solution to \textup{(A1)--(A3)} is unique:
$J(x) = \tfrac{1}{2}(x + x^{-1}) - 1$.
\end{theorem}

\begin{proof}[Proof sketch]
Set $H(t) = J_{\log}(t)+1$.  Axiom A2 transforms into the d'Alembert
equation $H(s+t)+H(s-t)=2H(s)H(t)$.  By the Acz\'el classification
theorem~\cite{Aczel1966}, the unique continuous solution with $H(0)=1$
and $H''(0)=1$ is $H(t)=\cosh t$, giving
$J(x)=\tfrac{1}{2}(x+x^{-1})-1$.
\end{proof}

\noindent The $\varphi$-ladder spacing $\ln\varphi$ distributes over the
8-tick recognition epoch ($2^D=8$, $D=3$ forced by $2^D=8$) giving the
per-epoch modulation $\delta_\varphi = \ln\varphi/2^D \approx 0.060$.
This is the same universal factor that appears in the gallium anomaly
resolution; it is a structural prediction of the $\varphi$-ladder and
8-tick period, not a free parameter.

%% ============================================================
\section{Introduction}
\label{sec:intro}
%% ============================================================

Big Bang Nucleosynthesis is one of the pillars of the hot Big Bang
model.  The predicted abundances of $^4$He, D, $^3$He, and $^7$Li
depend sensitively on the baryon-to-photon ratio $\eta$, independently
determined by Planck as $\eta = 6.1 \times 10^{-10}$.  For $^4$He and
D, BBN predictions agree with observations to better than $1\%$.  For
$^7$Li, they disagree by a factor of $\sim 3$.

The \emph{cosmological lithium problem}~\cite{Spite1982,Cyburt2016} is
this factor-of-3 discrepancy:
\begin{align}
  \left(\frac{^7\mathrm{Li}}{\mathrm{H}}\right)_{\mathrm{BBN}}
  &\approx (4.7 \pm 0.7) \times 10^{-10}, \\
  \left(\frac{^7\mathrm{Li}}{\mathrm{H}}\right)_{\mathrm{Spite}}
  &\approx (1.6 \pm 0.3) \times 10^{-10}.
\end{align}
The Spite plateau is the remarkably constant $^7$Li abundance observed
in metal-poor halo stars ([Fe/H] $< -1.5$), believed to reflect the
primordial value.

Proposed solutions include (i)~stellar depletion by mixing below the
convective zone, (ii)~non-standard BBN (extra neutrino species, varying
constants), and (iii)~nuclear physics uncertainties in reaction rates.
None fully resolves the problem without other tensions~\cite{Fields2011}.

Recognition Science~\cite{RS_Axioms} (RS) proposes a fourth option: the 8-tick recognition cycle suppresses
the ${}^3\mathrm{He}(\alpha,\gamma){}^7\mathrm{Be}$ reaction rate
during BBN by a derived fractional factor, accumulated over the
${}^7$Be production window.

%% ============================================================
\section{Standard BBN Lithium Production}
\label{sec:standard}
%% ============================================================

Primordial $^7$Li is produced primarily through:
\begin{equation}
  {}^3\mathrm{He} + {}^4\mathrm{He} \to {}^7\mathrm{Be} + \gamma,
  \qquad
  {}^7\mathrm{Be} + e^- \to {}^7\mathrm{Li} + \nu_e.
\end{equation}
The $^7$Be production window opens at $T \sim 0.4$~MeV (approximately
$100$~s after the Big Bang) and closes at $T \sim 0.1$~MeV
($\sim 1000$~s), as the temperature drops below the threshold for
nuclear reactions.  The $^7$Li yield is directly proportional to the
integrated ${}^3\mathrm{He}(\alpha,\gamma){}^7\mathrm{Be}$ rate over
this window.

%% ============================================================
\section{The RS Modulation Mechanism}
\label{sec:mechanism}
%% ============================================================

\subsection{Per-Cycle Rate Modulation}

\begin{theorem}[Rate Modulation Factor]
\label{thm:rate}
The 8-tick recognition cycle modulates nuclear reaction rates by a
fractional factor:
\begin{equation}
  \delta_\varphi = \frac{\ln\varphi}{2^D}
  = \frac{\ln\varphi}{8} \approx 0.060 \quad (6\%),
\end{equation}
where $D = 3$ is the number of spatial dimensions and $2^D = 8$ is
the 8-tick period.
\end{theorem}

\begin{proof}
The $\varphi$-ladder spacing in log-space is $\ln\varphi \approx 0.481$.
One 8-tick recognition epoch distributes this spacing uniformly across
the 8 ticks.  The per-tick fractional modulation is $\ln\varphi/8$.
For a nuclear reaction whose cross section depends on the $\varphi$-phase
of the recognition cycle, the per-epoch fractional modification to the
rate is:
\[
  \delta_\varphi = \frac{\ln\varphi}{2^D}
  = \frac{\ln\varphi}{8}
  = \frac{0.4812}{8}
  \approx 0.0601.
\]
The ${}^3\mathrm{He}(\alpha,\gamma){}^7\mathrm{Be}$ reaction involves a
transition between nuclear states of mass numbers $A = 3 + 4 = 7$
and $A = 7$, spanning one $\varphi$-rung in the nuclear binding energy
hierarchy.  The rung-crossing probability is modulated by
$\delta_\varphi$ per epoch.
\end{proof}

\begin{remark}
The same $6\%$ per-cycle modulation appears in the Gallium anomaly
(the $\nu$-nucleus cross section is suppressed by $57.7/74.0 \approx
0.780 \approx (1 - \delta_\varphi)^4$ over four recognition epochs,
corresponding to the $D+1 = 4$ independent ledger dimensions traversed
by a charged-current weak process).
The modulation factor $\delta_\varphi = \ln\varphi/8$ is thus a
universal RS prediction for electromagnetic and weak nuclear processes
that cross $\varphi$-rungs.
\end{remark}

\subsection{Cumulative Suppression}

\begin{theorem}[Cumulative Suppression Factor]
\label{thm:cumulative}
The cumulative suppression factor for $^7$Be production over $N$
effective recognition epochs is:
\begin{equation}
  F_{\rm sup}(N) = (1 - \delta_\varphi)^N = \left(1 - \frac{\ln\varphi}{8}\right)^N.
\end{equation}
For $N \approx 18$:
\begin{equation}
  F_{\rm sup}(18) = (1 - 0.0601)^{18} \approx 0.328 \approx \frac{1}{3}.
\end{equation}
\end{theorem}

\begin{proof}
At each effective recognition epoch, the $^7$Be production rate is
suppressed by factor $(1 - \delta_\varphi)$.  After $N$ independent
epochs, the rate is $(1 - \delta_\varphi)^N$.  Numerically:
\[
  (1 - 0.0601)^{18} = (0.9399)^{18}.
\]
Computing $\ln(0.9399^{18}) = 18 \times \ln(0.9399) = 18 \times
(-0.0619) = -1.114$, so $(0.9399)^{18} = e^{-1.114} \approx 0.328$.
\end{proof}

\begin{remark}[Status of $N \approx 18$]
\label{rem:N}
The value $N \approx 18$ is derived by \emph{inverting} the
required suppression factor---it is the unique integer satisfying
$F_{\rm sup}(N) \approx 1/3$ given $\delta_\varphi \approx 0.060$.
Explicitly:
\[
  N = \frac{\ln(1/3)}{\ln(1 - \delta_\varphi)} = \frac{-1.099}{-0.0619} \approx 17.8.
\]
This is not independently derived from first principles.  Establishing
$N = 18$ from the RS Hamiltonian would require:
\begin{enumerate}[nosep,label=(\roman*)]
  \item Identifying which nuclear reaction eigenstates couple to the
  recognition ledger during BBN;
  \item Integrating the modified Boltzmann network over the ${}^7$Be
  production window ($100$~s--$1000$~s); and
  \item Counting the number of $\varphi$-rung crossings in the
  ${}^3$He--$\alpha$ scattering matrix element.
\end{enumerate}
An order-of-magnitude consistency check: the nuclear matrix element
for ${}^3\mathrm{He}(\alpha,\gamma){}^7\mathrm{Be}$ involves
$A = 7$ binding relative to $A = 3 + 4$ scattering states, which
span $\sim 2$--$3$ $\varphi$-rungs in the nuclear binding energy
ladder ($E_{\rm bind} \sim 1$--$8$~MeV maps to rungs $\sim 38$--$41$).
The integrated effective epochs for $\sim 20$--$30$ rung transitions
over the reaction network is consistent with $N \sim 18$ at the
order-of-magnitude level, but a precise first-principles derivation
remains to be completed.
\end{remark}

\subsection{Predicted Lithium Abundance}

\begin{corollary}[RS-Corrected Lithium Abundance]
\label{cor:lithium}
The RS-predicted primordial $^7$Li abundance is:
\begin{equation}
  \left(\frac{^7\mathrm{Li}}{\mathrm{H}}\right)_{\mathrm{RS}}
  = F_{\rm sup}(18) \times \left(\frac{^7\mathrm{Li}}{\mathrm{H}}\right)_{\mathrm{BBN}}
  = 0.328 \times 4.7 \times 10^{-10}
  \approx 1.54 \times 10^{-10},
\end{equation}
consistent with the Spite plateau value $(1.6 \pm 0.3) \times
10^{-10}$ within $0.2\sigma$.
\end{corollary}

%% ============================================================
\section{Preservation of Other Abundances}
\label{sec:preserve}
%% ============================================================

\begin{proposition}[$^4$He Insensitivity]
\label{prop:helium}
The $^4$He mass fraction $Y_p$ is insensitive to the $\varphi$-modulation
of the ${}^3\mathrm{He}(\alpha,\gamma){}^7\mathrm{Be}$ rate.  The
RS-corrected value is $Y_p \approx 0.247$, consistent with the observed
$Y_p = 0.245 \pm 0.003$~\cite{Aver2015}.
\end{proposition}

\begin{proof}[Physical argument]
The $^4$He mass fraction is determined by the neutron-to-proton ratio
$n/p$ at weak-interaction freeze-out ($T \sim 0.7$~MeV, $t \sim 1$~s).
Subsequent $^4$He synthesis proceeds via
$p + n \to d + \gamma$ and $d + d \to {}^4\mathrm{He}$, which
completes rapidly and converts virtually all available neutrons.

The ${}^3\mathrm{He}(\alpha,\gamma){}^7\mathrm{Be}$ reaction opens
only at $T \lesssim 0.4$~MeV ($t \gtrsim 100$~s), \emph{after}
$^4$He synthesis is complete.  Therefore the $\varphi$-modulation of
this rate has no effect on $Y_p$.  The $^4$He yield depends on $n/p$
at freeze-out and the efficiency of $^4$He synthesis, neither of
which involves the ${}^3\mathrm{He}(\alpha,\gamma){}^7\mathrm{Be}$
channel.

The weak interaction rates governing $n/p$ freeze-out are suppressed
by separate $\varphi$-modulation, but $Y_p$ is logarithmically
insensitive to them: a $6\%$ change in the freeze-out rate shifts
$n/p$ by $\lesssim 0.3\%$ and $Y_p$ by $\lesssim 0.001$, within
the observational uncertainty.
\end{proof}

\begin{proposition}[Deuterium Preservation]
\label{prop:deuterium}
The deuterium ratio D/H is preserved at $\mathrm{D/H} \approx
2.5 \times 10^{-5}$, consistent with the observed
$(2.53 \pm 0.04) \times 10^{-5}$~\cite{Cooke2018}.
\end{proposition}

\begin{proof}[Physical argument]
Deuterium production and destruction both proceed via strong and
electromagnetic reactions at $T \sim 0.07$~MeV.  The equilibrium D/H
is set by the balance between $p + n \to d + \gamma$ (forward) and
$d + \gamma \to p + n$ (reverse), plus destruction by
$d + p \to {}^3\mathrm{He} + \gamma$, etc.  In RS, both forward and
reverse reaction rates receive the same $\delta_\varphi$ modulation per
epoch, since they involve the same $\varphi$-rung structure.  To first
order in $\delta_\varphi$, the equilibrium D/H is unchanged; the leading
correction is of order $\delta_\varphi^2 \approx 0.4\%$, far below the
$\sim 1.6\%$ observational uncertainty on D/H.
\end{proof}

%% ============================================================
\section{Experimental Comparison}
\label{sec:compare}
%% ============================================================

\begin{center}
\renewcommand{\arraystretch}{1.3}
\begin{tabular}{@{}lccc@{}}
\toprule
\textbf{Species} & \textbf{Standard BBN} & \textbf{RS-Corrected} &
\textbf{Observed} \\
\midrule
$^7$Li/H ($\times 10^{-10}$) & $4.7 \pm 0.7$ & $\sim 1.54$ &
$1.6 \pm 0.3$ \\
$Y_p$ ($^4$He) & $0.247$ & $0.247$ & $0.245 \pm 0.003$ \\
D/H ($\times 10^{-5}$) & $2.5$ & $2.5$ & $2.53 \pm 0.04$ \\
\bottomrule
\end{tabular}
\end{center}

%% ============================================================
\section{Predictions and Falsification}
\label{sec:predict}
%% ============================================================

\begin{enumerate}
  \item \textbf{Laboratory rates unchanged.}  The $\varphi$-modulation
  is a cosmological effect driven by the density and temperature
  evolution during BBN.  Laboratory measurements of the
  ${}^3\mathrm{He}(\alpha,\gamma){}^7\mathrm{Be}$ cross section at
  low energies should \emph{not} resolve the lithium problem,
  since the mechanism is in the cosmological context, not in the
  nuclear rate itself.  Continued agreement between laboratory rates
  and the best current evaluations~\cite{Cyburt2016} is predicted.

  \item \textbf{$^6$Li/H.}  The same $\delta_\varphi$ modulation
  applies to $^6$Li production via ${}^4\mathrm{He}(d,\gamma){}^6\mathrm{Li}$.
  The RS-predicted $^6$Li/$^7$Li ratio and primordial $^6$Li abundance
  can be computed from Theorem~\ref{thm:cumulative} once $N^{(6)}$
  (the effective epochs for $^6$Li production) is derived.

  \item \textbf{Stellar depletion not required.}  RS predicts that
  the Spite plateau reflects the true primordial $^7$Li abundance,
  with no need for stellar depletion.  Observational evidence for
  lithium depletion by mixing (e.g., a slope in $^7$Li vs.\ effective
  temperature in warm metal-poor stars) would be inconsistent with
  this prediction.

  \item \textbf{Falsifiers.}
  \begin{itemize}[nosep]
    \item If a first-principles RS calculation of $N$ yields
    $N \notin [15, 22]$, the quantitative agreement is accidental.
    \item If new observations revise the Spite plateau to
    $^7$Li/H $> 2.5 \times 10^{-10}$ or $< 1.0 \times 10^{-10}$,
    the RS mechanism is falsified.
    \item If stellar models successfully explain $^7$Li depletion from
    the standard BBN value down to the Spite plateau, the RS
    mechanism becomes unnecessary.
  \end{itemize}
\end{enumerate}

%% ============================================================
\section{Conclusion}
\label{sec:conclusion}
%% ============================================================

Recognition Science provides a specific mechanism for the cosmological
lithium problem: the $\varphi$-modulation of the
${}^3\mathrm{He}(\alpha,\gamma){}^7\mathrm{Be}$ rate by
$\delta_\varphi = \ln\varphi/2^D \approx 6\%$ per recognition epoch.
The modulation factor is derived from the RS $\varphi$-ladder and
8-tick period.  After $N \approx 18$ effective epochs---consistent
with the order-of-magnitude count of $\varphi$-rung crossings during
the BBN ${}^7$Be production window---the predicted $^7$Li/H $\approx
1.54 \times 10^{-10}$ matches the Spite plateau $(1.6 \pm 0.3)
\times 10^{-10}$.  The $^4$He and D abundances are insensitive to
this channel-specific suppression and remain in agreement with
observation.  The key outstanding task is to derive $N = 18$
from the RS modified Boltzmann network for the BBN reaction chain,
confirming that the quantitative match is not coincidental.

%% ============================================================
\begin{thebibliography}{99}

\bibitem{Cyburt2016}
R.~H.~Cyburt et al., Rev.\ Mod.\ Phys.\ \textbf{88}, 015004 (2016).

\bibitem{Fields2011}
B.~D.~Fields, Ann.\ Rev.\ Nucl.\ Part.\ Sci.\ \textbf{61}, 47 (2011).

\bibitem{Spite1982}
F.~Spite and M.~Spite, A\&A \textbf{115}, 357 (1982).

\bibitem{Aver2015}
E.~Aver, K.~A.~Olive, and E.~D.~Skillman, JCAP \textbf{1507}, 011
(2015).

\bibitem{Cooke2018}
R.~J.~Cooke, M.~Pettini, and C.~C.~Steidel, ApJ \textbf{855}, 102
(2018).

\bibitem{RS_Axioms}
J.~Washburn, ``The Algebra of Reality: A Recognition Science Derivation
of Physical Law,'' \emph{Axioms} \textbf{15}(2), 90 (2026).
\texttt{doi:10.3390/axioms15020090}

\bibitem{Aczel1966}
J.~Acz\'el, \emph{Lectures on Functional Equations and Their
Applications} (Academic Press, 1966).

\end{thebibliography}

\end{document}
